%%%%%%%%%%%%%%%%%%%%%%%%%%%%%%%%%%%%%%%%%%%%%%%%%%%%%%%%%%%%%%%%%%%%%%%%%%%%%%%%
\section{Introduction}
% 背景意义：复杂3D穿越不仅需要落脚信息，也需要身体可通行空间信息。
Humanoid robots are expected to move through human-scale environments that contain stairs, cluttered passages, protruding structures and other three-dimensional constraints. For successful traversal, the controller depends on not only estimating where the feet can step on, but also understanding whether the whole body can safely pass through the surrounding space, including overhead and lateral obstacles that may cause collision or loss of balance. This makes perceptive traversal a coupled problem of terrain understanding, body-level clearance reasoning and real-time whole-body control.

% 一方面：现有表示在完整3D语义和实时计算之间难以兼顾。
However, onboard perception for legged robots is constrained by strict latency and bandwidth requirements. Existing representations therefore tend to emphasize either efficiency or geometric completeness. Conventional 2.5D elevation maps~\cite{Fankhauser2014RobotCentricElevationMapping,Fankhauser2018ProbabilisticTerrainMapping,ElevationMap_GPU_2022,long2025learning,wang2025beamdojo,ren2025vb,he2025attention,videomimic,Multi-Layer_ElevationMap_2025,Roth_2025_LearnedPerceptiveForwardDynamicsModel,Zhang_2026_AME2} provide compact terrain geometry for foothold reasoning, but they are intrinsically blind to overhangs, side walls and body-level clearance constraints. Depth-based representations~\cite{Li_2024_MOVE,zhuang2024humanoidRobot,zhuang2024humanoid,Yu_2025_START,Castillo_2025_LearningTerrainAwareBipedalLocomotion,Solano_2025_StandWalkNavigate,Sun_2025_DPL,Wang_2025_LocoMamba,zhu2026hiking} offer direct visual observations, yet they can be affected by view changes, occlusion and illumination variation. Volumetric approaches~\cite{Ben_2025_Gallant,Niijima_2025_Real-timeMulti-PlaneSegmentation,Wang_2025_Omni-Perception} encode richer 3D geometry, but their dimensionality and processing cost make high-frequency policy learning and deployment difficult. These limitations necessitate a representation that preserves 3D traversability while remaining computationally compact~\cite{2020arXiv200716100T,2018arXiv181205784L,2022arXiv220108117M}.

% 另一方面：固定感知尺度无法适应速度和局部复杂度的变化。
A second limitation is that perception is usually acquired with a fixed spatial range~\cite{Ben_2025_Gallant,Multi-Layer_ElevationMap_2025,long2025learning}, which is misaligned with the changing information demand of locomotion: fast motion needs a longer horizon while cluttered traversal needs higher local resolution. A perceptive controller should therefore decide not only how to act given perception, but also how much spatial context should be sensed.

% 为了解决这两个问题：我们提出双投影，自适应半径解决感知尺度问题。
To address these two issues, this paper proposes \textbf{ADAPT} (\textbf{A}daptive \textbf{D}ual-projection \textbf{A}rchitecture for \textbf{P}erceptive \textbf{T}raversal) that restructures local 3D geometry into two orthogonal and information-dense projections. A horizontal elevation map describes ground geometry for stable foothold selection, while a vertical distance map describes body-level traversable space for avoiding overhangs and obstacles. Instead of using a fixed perceptual horizon, ADAPT further exposes the sensing radius of the vertical distance map to the reinforcement learning policy as a control variable. The policy can thus expand its perceptual horizon during fast traversal and contract it for fine-grained local reasoning in cluttered scenes. In this way, ADAPT jointly addresses what spatial information should be represented and how much of it should be sensed at each moment.

% 实验意义：实验指标对应鲁棒性、效率和真机可部署性。
Empirical evaluation demonstrates that this design improves both robustness and efficiency. ADAPT achieves a $94.7\%$ success rate across a series of 3D negotiation tasks and surpasses baselines with fixed sensing ranges by $30\%$. The dual-projection architecture reduces observation dimensionality by $98.6\%$ and decreases the training time per iteration to $32.7\%$ of that required by voxel-based alternatives. The learned adaptive sensing mechanism further reduces computational overhead by $56\%$ by matching resource consumption to motion complexity. These results indicate that ADAPT is not only a compact perception representation, but also a perception-control coupling mechanism suitable for real-time humanoid traversal.

The main contributions of this work are:
\begin{itemize}
    \item Proposing a highly efficient formulation that decouples complex 3D environments into compact projections. This paradigm reduces observational dimensionality while rigorously preserving essential geometric constraints, thereby ensuring the high-frequency perception strictly required for agile locomotion.
    \item Introducing a bidirectional closed-loop mechanism that embeds perceptual configurations directly into the reinforcement learning action space. This formulation empowers the robot to autonomously modulate its sensory settings based on instantaneous dynamic states and environmental complexity, thereby achieving an optimal trade-off between navigational safety and precise spatial perception.
    \item Demonstrating robust zero-shot sim-to-real transfer on a Unitree G1 humanoid robot with successful navigation through diverse structural challenges.
\end{itemize}
